\documentclass[11pt,a4paper]{article}
\usepackage[T1]{fontenc}
\usepackage[utf8]{inputenc}
% O pacote de hifenização em português não está presente em todas as instalações
% mínimas do TeX Live. O texto continua em UTF-8 e compila sem essa dependência.
\usepackage{lmodern}
\usepackage[a4paper,left=2.3cm,right=2.1cm,top=2.2cm,bottom=2.1cm]{geometry}
\usepackage{amsmath,booktabs,longtable,tabularx,array,multirow}
\usepackage{graphicx,xcolor,microtype,enumitem,ragged2e,float,pdflscape}
\usepackage{caption,fancyhdr,titlesec}
\usepackage[hidelinks]{hyperref}
\renewcommand{\contentsname}{Sumário}
\renewcommand{\refname}{Referências}
\renewcommand{\tablename}{Tabela}

\definecolor{azul}{HTML}{2A78D6}
\definecolor{laranja}{HTML}{EB6834}
\definecolor{tinta}{HTML}{14181F}
\definecolor{tinta2}{HTML}{525B66}
\definecolor{linha}{HTML}{DFE3E8}
\definecolor{fundo}{HTML}{F4F5F7}
\captionsetup{font=small,labelfont={bf,color=tinta},textfont={color=tinta2},
  justification=justified,singlelinecheck=false,skip=5pt}
\titlespacing*{\section}{0pt}{12pt plus 2pt minus 2pt}{5pt}
\titlespacing*{\subsection}{0pt}{9pt plus 2pt minus 2pt}{3pt}
\setlength{\parskip}{2pt plus 1pt}
\setlist{topsep=3pt,itemsep=2pt,parsep=1pt}
\newcolumntype{Y}{>{\RaggedRight\arraybackslash}X}
\newcommand{\destaque}[2]{\begin{center}\begin{minipage}{.96\textwidth}
  \colorbox{fundo}{\begin{minipage}{\dimexpr\linewidth-2\fboxsep}
  \textbf{\color{azul}#1}\par\small #2\end{minipage}}\end{minipage}\end{center}}

% Preencher apenas com informação confirmada pelo grupo.
\newcommand{\numerogrupo}{3}

\pagestyle{fancy}\fancyhf{}
\fancyhead[L]{\small\color{tinta2}Cloud Computing e Cloud Storage --- AWS}
\fancyfoot[C]{\small\thepage}
\renewcommand{\headrulewidth}{.4pt}
\hypersetup{pdftitle={Cloud Computing e Cloud Storage: estudo de caso AWS},
  pdfauthor={Grupo de Construcao de Banco de Dados --- UFRJ}}

\begin{document}
\thispagestyle{empty}
\begin{center}
{\large\scshape Universidade Federal do Rio de Janeiro}\\[2pt]
{\normalsize Instituto de Matemática --- Departamento de Ciência da Computação}\\[2pt]
{\normalsize Disciplina: Construção de Banco de Dados}\\[2pt]
{\normalsize Professor: Milton Ramirez}\\[17mm]
{\color{azul}\rule{\textwidth}{1.2pt}}\\[6mm]
{\LARGE\bfseries Cloud Computing e Cloud Storage}\\[4mm]
{\Large Estudo de caso: Amazon Web Services (AWS)}\\[4mm]
{\large Arquitetura, serviços, segurança e análise de 3,5 PB}\\[6mm]
{\color{azul}\rule{\textwidth}{1.2pt}}\\[8mm]
{\large\bfseries Grupo \numerogrupo}\\[6mm]
\begin{tabular}{@{}ll@{}}\toprule
\textbf{Nome completo} & \textbf{DRE}\\\midrule
Bernardo Brandão Pozzato Carvalho Costa & 123289593\\
Enzo de Carvalho Sampaio & 123386206\\
Gabriel Schmitz Corrêa Rizawinsk & 123225573\\
Guilherme En Shih Hu & 123224674\\
Raphael Henrique da Silva Pereira & 123311073\\
Vivian Maria da Silva e Souza & 123205793\\\bottomrule
\end{tabular}\\[12mm]
{\normalsize Rio de Janeiro --- Setembro de 2026}
\end{center}
\clearpage

\section*{Resumo}
\addcontentsline{toc}{section}{Resumo}
Este trabalho explica computação e armazenamento em nuvem e analisa a AWS como provedor. Partimos
das cinco características do NIST, distinguimos IaaS, PaaS, SaaS, DBaaS e DWaaS e comparamos
nuvens pública, privada e híbrida. Em seguida relacionamos requisitos de banco de dados aos
serviços EC2, S3, EBS, classes S3 Glacier, RDS, Aurora, DynamoDB e Redshift. O estudo econômico
trata 3,5~PB como um problema de arquitetura, não como simples multiplicação de preço por GB:
separa armazenamento, processamento, requisições, recuperação, transferência, proteção e suporte.
Com preços públicos da região us-east-1 consultados em 6 de setembro de 2026, 3,5~PB decimais
correspondem a aproximadamente 3.259.629 unidades de cobrança binárias da AWS. Custariam cerca de
US\$~69.015/mês somente em S3 Standard, US\$~40.745 em Standard-IA,
US\$~11.735 em Glacier Flexible Retrieval ou US\$~3.227 em Deep Archive. Esses valores não são
alternativas equivalentes: possuem latências, cobranças de recuperação e durações mínimas distintas.
Propomos, por isso, um data lake em S3 com dados colunares e particionados, catálogo, políticas de
ciclo de vida e computação escolhida segundo o perfil da consulta. O Post-Mortem registra os
prompts, as intervenções humanas ainda pendentes de confirmação e as correções de erros da IA.

\noindent\textbf{Palavras-chave:} computação em nuvem; AWS; armazenamento de objetos; DBaaS;
DWaaS; data lake; custo total de propriedade; 3,5 PB.
\clearpage
\tableofcontents
\clearpage

\section*{Introdução}
\addcontentsline{toc}{section}{Introdução}
A nuvem troca a compra antecipada de infraestrutura por recursos medidos, provisionados por API e
compartilhados com isolamento. Para um SBD, essa troca altera o local do controle: o cliente pode
administrar sistema operacional e SGBD em EC2, delegar tarefas operacionais ao RDS ou Aurora, ou
usar serviços com modelos de dados próprios, como DynamoDB e Redshift. O ganho de elasticidade não
elimina decisões de engenharia. Região, zonas, consistência, latência, modelo de falha, RPO, RTO,
governança e custo continuam sob responsabilidade do arquiteto.

O enunciado pede também uma decisão para 3,5~PB. Sem caracterizar o conjunto de dados, não existe
um único preço tecnicamente honesto: uma base transacional, um data warehouse e um arquivo regulatório
exigem serviços diferentes. Este relatório explicita premissas e oferece três cenários comparáveis.

Nossa tese é que, nessa escala, a primeira decisão não é escolher uma instância: é separar a
persistência durável da capacidade de computação. Objetos em S3 formam a camada de verdade para o
acervo analítico; engines diferentes processam subconjuntos conforme latência e concorrência; bancos
operacionais preservam apenas o estado que realmente exige transações. Essa separação evita usar
EBS, RDS ou Aurora como se fossem recipientes indiferenciados de petabytes.

\begin{table}[H]\centering\caption{Checklist do enunciado.}\label{tab:checklist}
\small\begin{tabularx}{\textwidth}{@{}Yl@{}}\toprule
\textbf{Exigência} & \textbf{Localização}\\\midrule
Foundations and deployment: NIST, características e implantação & Seção~\ref{sec:foundations}\\
Service Models and Evaluation: IaaS, PaaS, SaaS e critérios & Seção~\ref{sec:service-evaluation}\\
EC2; S3, EBS e Glacier; RDS, Aurora, DynamoDB e Redshift & Seção~\ref{sec:aws}\\
Custo, vantagens e desvantagens para 3,5 PB & Seção~\ref{sec:custo}\\
Post-Mortem, prompts, autoria e correções da IA & Seção~\ref{sec:postmortem}\\
Folha de rosto com nomes e DRE & página 1\\\bottomrule
\end{tabularx}\end{table}

\subsection*{Método e hierarquia das fontes}
\addcontentsline{toc}{subsection}{Método e hierarquia das fontes}
O estudo combina os livros indicados no enunciado com normas e documentação operacional. Usamos o
NIST para a definição de nuvem; livros-texto para modelos e implicações de SBD; documentação AWS
para semântica, limites e preços. Afirmações mutáveis registram região e data. Preço publicado é
tratado como estimativa, não proposta comercial. Quando a documentação usa ``GB'' para uma unidade
binária, preservamos a regra de cobrança e mostramos a conversão em vez de misturar PB, PiB, GB e
GiB \cite{s3pricing}.

As comparações seguem quatro perguntas: (1) qual abstração o consumidor recebe; (2) quem mantém cada
camada; (3) qual SLO pode ser construído; e (4) quais variáveis compõem o custo total. Isso aproxima
o relatório do problema de banco de dados: capacidade sem semântica de escrita, durabilidade sem
recuperação e preço sem workload não permitem uma decisão.

\section{Bases Conceituais}\label{sec:foundations}
Esta seção estabelece a base conceitual para as decisões posteriores. Primeiro, apresenta a definição
do NIST e suas características essenciais; depois, compara os modelos público, privado e híbrido.
Esses modelos descrevem a governança e a integração da infraestrutura, e não devem ser confundidos
com os modelos de serviço, que delimitam responsabilidades operacionais.

\subsection{Definição de computação em nuvem segundo o NIST}
O NIST define computação em nuvem como um modelo de acesso conveniente e sob demanda, pela rede,
a um conjunto compartilhado de recursos configuráveis que podem ser rapidamente provisionados e
liberados com pouco esforço de gestão \cite{nist}. Em outras palavras, a organização deixa de
tratar servidores, redes e armazenamento apenas como ativos físicos comprados antecipadamente e passa
a consumi-los com recursos que podem ser solicitados, configurados e medidos conforme a necessidade.

Para diferenciar esse modelo de outras formas de hospedar sistemas, o NIST estabelece cinco
características essenciais: autoatendimento sob demanda, amplo acesso por rede, agrupamento de
recursos, elasticidade rápida e serviço mensurável. Elas descrevem tanto a experiência do consumidor
quanto a forma como o provedor organiza sua infraestrutura. A virtualização é uma tecnologia que
frequentemente viabiliza essa organização, pois permite criar ambientes isolados sobre o mesmo
hardware. Contudo, ela não é sinônimo de nuvem: uma máquina virtual criada manualmente em um único
servidor pode ser virtualizada, mas não oferece, por si só, os cinco critérios do NIST.

\subsection{As cinco características essenciais}
A primeira característica, o \textbf{autoatendimento sob demanda}, permite que o próprio consumidor
crie, redimensione ou encerre um recurso por console, API ou código de infraestrutura, sem depender
de uma solicitação manual ao operador. O \textbf{amplo acesso por rede} torna esses recursos
disponíveis por protocolos padronizados a clientes diferentes, como aplicações web, notebooks e
servidores. Isso não significa que o dado esteja aberto na Internet: autenticação, redes privadas e
políticas continuam determinando quem pode acessar o serviço e por qual caminho.

O \textbf{agrupamento de recursos} explica como o provedor atende muitos consumidores. A
virtualização abstrai CPU, memória, rede e dispositivos físicos em recursos lógicos isolados; o
\emph{pooling} administra a capacidade de muitos equipamentos como um conjunto, alocando-a conforme
a demanda sem que cada cliente escolha o servidor físico exato. Essa modalidade é chamada de
\emph{multi-tenancy}: a infraestrutura é compartilhada, mas identidade, rede, execução e dados
precisam permanecer isolados. Assim, o risco não desaparece com a abstração; ele passa a exigir
controles lógicos, auditoria e evidências de que o isolamento é efetivo.

As duas características restantes relacionam capacidade e custo. \textbf{Elasticidade rápida} é a
possibilidade de aumentar ou reduzir recursos em tempo compatível com a variação da demanda. Ela não
deve ser confundida com escalabilidade: escala vertical aumenta CPU ou memória de uma única máquina,
enquanto escala horizontal adiciona máquinas; elasticidade é a capacidade de executar esse ajuste de
modo ágil, idealmente automático. Um sistema pode ser escalável, mas não elástico, se a expansão
exigir uma janela manual de manutenção. Em sistemas de banco de dados, a escala horizontal é mais
difícil porque escrita e estado exigem particionamento, coordenação, consistência e redistribuição de
dados. Por fim, o \textbf{serviço mensurável} registra o consumo e o cobra em unidades como horas
de computação, GB-mês ou número de requisições; portanto, a arquitetura também influencia a conta
financeira.

\begin{table}[H]\centering\caption{Das características NIST às consequências de arquitetura.}
\small\begin{tabularx}{\textwidth}{@{}p{3.1cm}Y Y@{}}\toprule
\textbf{Característica} & \textbf{Mecanismo observável} & \textbf{Implicação para SBD}\\\midrule
Autoatendimento & API/IaC cria e remove recursos sem chamado & velocidade exige guardrails, quotas e revisão de configuração\\
Amplo acesso por rede & endpoints e protocolos padronizados & latência, identidade e caminho privado entram no projeto\\
Pool de recursos & alocação dinâmica e isolamento entre clientes & localização física é abstraída; região e AZ continuam relevantes\\
Elasticidade rápida & capacidade cresce/reduz conforme demanda & conexões, réplicas e partições precisam acompanhar o compute\\
Serviço mensurável & consumo metrado por dimensões do serviço & formato físico e padrão de acesso tornam-se variáveis financeiras\\\bottomrule
\end{tabularx}\end{table}

\subsection{Cloud Storage}
O armazenamento em nuvem (\emph{Cloud Storage}) consiste na oferta de persistência de dados como um serviço gerenciado. Ele se divide fundamentalmente em três categorias arquiteturais, cada uma com uma forma diferente de entregar e acessar os dados:

\begin{enumerate}
    \item \textbf{Armazenamento de Blocos:} Funciona como um "disco virtual" bruto que é conectado a um servidor (\emph{host}), permitindo que o sistema operacional crie nele seu próprio sistema de arquivos. Na AWS, o exemplo principal é o \textbf{EBS}.
    \item \textbf{Armazenamento de Arquivos:} Disponibiliza uma estrutura tradicional de pastas e diretórios compartilhados por meio de protocolos de rede (como NFS ou SMB). Na AWS, o exemplo é o \textbf{EFS}.
    \item \textbf{Armazenamento de Objetos:} Organiza os dados em repositórios chamados \emph{buckets} e gerencia cada arquivo como um "objeto" isolado, acessado por meio de chaves, metadados e chamadas de API (via HTTP). Na AWS, o principal exemplo é o \textbf{S3}. As classes do \textbf{S3 Glacier} são apenas opções de arquivamento de longo prazo e baixo custo que ficam \emph{dentro} do próprio ecossistema do S3; elas não funcionam como discos de servidor que podem ser montados diretamente.
\end{enumerate}

A diferença entre essas categorias não está apenas na velocidade, mas na \textbf{forma como os dados são lidos e gravados}:

\begin{itemize}[leftmargin=*]
    \item \textbf{Modelo de Blocos:} entrega endereços de disco brutos e deixa o controle de arquivos e concorrência por conta do servidor. Bancos de dados relacionais tradicionais dependem de gravações frequentes e alterações em pequenas partes de arquivos (páginas de dados), o que se encaixa perfeitamente nesse modelo.
    \item \textbf{Modelo de Objetos:} trata o arquivo como um bloco único e indivisível. Se for preciso alterar apenas uma linha de um arquivo no S3, o sistema geralmente precisa substituir ou versionar o objeto inteiro por meio da API. Por isso, os \emph{data lakes} utilizam o S3 para armazenar arquivos grandes e \emph{imutáveis} (que não mudam após criados), otimizados para análises em lote.
\end{itemize}

Em suma, tentar forçar o S3 a se comportar como um disco rígido tradicional (EBS) por meio de camadas de compatibilidade não altera a sua natureza fundamental, gerando problemas de desempenho e arquitetura.
\subsection{Consistência, durabilidade e disponibilidade}

Em sistemas de armazenamento em nuvem e arquiteturas de banco de dados, os conceitos de consistência, durabilidade e disponibilidade representam eixos independentes, porém complementares, que definem o comportamento e as garantias operacionais de uma plataforma.

\begin{itemize}[leftmargin=*]
    \item \textbf{Consistência:} Responde à pergunta \textit{``qual versão de um dado a leitura observa imediatamente após uma operação de escrita?''}. O Amazon S3, por exemplo, oferece consistência forte de leitura após operações de \texttt{PUT} e \texttt{DELETE} em objetos novos ou existentes \cite{s3consistency}. Contudo, consistência técnica ao nível de infraestrutura não valida a semântica da aplicação: o sistema pode gravar com sucesso um estado logicamente inválido ou permitir a exclusão acidental de um objeto correto.
    \item \textbf{Durabilidade:} Responde à pergunta \textit{``qual é a probabilidade estatística de que o dado aceito e armazenado permaneça intacto ao longo do tempo, sem corrupção ou perda?''}. É uma métrica voltada para a resistência física e lógica contra falhas catastróficas de hardware ou de infraestrutura de armazenamento.
    \item \textbf{Disponibilidade:} Responde à pergunta \textit{``qual é a probabilidade de que uma operação receba uma resposta bem-sucedida no momento em que é solicitada?''}. Trata-se da operacionalidade contínua do serviço sob a perspectiva do cliente, mensurada por meio de acordos de nível de serviço (SLA) e indicadores de desempenho (\emph{SLIs}).
\end{itemize}

Nenhum mecanismo de proteção isolado é capaz de mitigar todas as classes de risco inerentes a um ambiente de grande escala. Cada camada de segurança e redundância endereça uma vulnerabilidade distinta:
\begin{enumerate}
    \item \textit{A replicação física interna do provedor} protege contra falhas de hardware e indisponibilidades de componentes de infraestrutura dentro do modelo do serviço.
    \item \textit{O versionamento} protege contra a sobrescrita indesejada ou a exclusão lógica equivocada de dados.
    \item \textit{O backup independente e isolado} mitiga perdas decorrentes de comprometimento de credenciais ou falhas no domínio administrativo principal.
    \item \textit{A replicação entre regiões geográficas} serve como barreira contra desastres de grande escala que afetem uma área inteira.
\end{enumerate}

Portanto, o projeto de arquitetura exige uma definição clara de ameaças, objetivos de ponto de recuperação (\emph{RPO}), objetivos de tempo de recuperação (\emph{RTO}) e o rigoroso controle de privilégios sobre quem detém a permissão para apagar cópias de recuperação.

\destaque{Distinção central}{Durabilidade mede a probabilidade de o dado persistir; disponibilidade
mede a probabilidade de o serviço responder; consistência descreve o que uma leitura pode observar.
Uma porcentagem de ``11 noves'' de durabilidade do S3 (99,999999999\%) não promete 100\% de disponibilidade,
nem substitui a necessidade de versionamento, cópias redundantes entre regiões e testes periódicos de restauração.}

\section{Modelos de implantação}

A escolha de um modelo de implantação determina onde a infraestrutura física reside, quem a opera e como o acesso e a governança são estruturados. Essa classificação difere dos modelos de serviço (como IaaS, PaaS ou SaaS), pois descreve a arquitetura de propriedade e isolamento organizacional, e não o nível de abstração da plataforma.

\subsection{Pública}

Na \textbf{nuvem pública}, um provedor terceirizado (como AWS, Azure ou GCP) aloca recursos computacionais sob demanda, compartilhados entre múltiplos clientes por meio de virtualização e isolamento lógico (\emph{multi-tenancy}). 

\begin{itemize}[leftmargin=*]
    \item \textbf{Vantagens:} Elasticidade quase ilimitada, eliminação de CAPEX (despesa de capital) em hardware físico e acesso imediato a um vasto ecossistema de serviços gerenciados.
    \item \textbf{Mitigação de riscos de rede:} ``Pública'' refere-se à propriedade da infraestrutura e ao modelo de entrega comercial, e não à exposição desprotegida na Internet. O uso de redes virtuais isoladas (VPCs), túneis VPN, conexões dedicadas (como AWS Direct Connect) e \emph{endpoints} privados assegura que o tráfego corporativo ocorra em canais controlados.
    \item \textbf{Responsabilidade do cliente:} Embora o provedor gerencie o datacenter e os hosts físicos, a organização continua inteiramente responsável pela configuração de identidades (IAM), políticas de acesso, criptografia, monitoramento de quotas e controle de custos (FinOps).
\end{itemize}

\subsection{Privada}

Na \textbf{nuvem privada}, a infraestrutura é dedicada de forma exclusiva a uma única organização, podendo estar localizada fisicamente em seu próprio datacenter (\emph{on-premises}) ou operada por um fornecedor terceirizado em um ambiente hospedado.

\begin{itemize}[leftmargin=*]
    \item \textbf{Vantagens:} Máximo controle sobre a topologia de rede, segurança física direta, customização avançada de hardware e atendimento rigoroso a mandatos de conformidade ou residência de dados.
    \item \textbf{Limitações e custos:} A capacidade é estática e finita, exigindo planejamento financeiro antecipado e investimento de capital. A organização absorve integralmente os custos de atualização tecnológica, redundância de energia, refrigeração e manutenção de equipe especializada.
    \item \textbf{Segurança operacional:} Uma nuvem privada não é inerentemente mais segura que a pública; a segurança depende estritamente da rigidez dos controles físicos, segmentações lógicas e processos de auditoria aplicados internamente pela equipe da instituição.
\end{itemize}

\subsection{Híbrida e Padrões de Integração}

A \textbf{nuvem híbrida} integra ambientes de nuvem privada (\emph{on-premises}) com recursos de nuvem pública em uma única arquitetura coesa. A mera existência de servidores locais e instâncias na nuvem não caracteriza um modelo híbrido; é exigida a interoperabilidade contínua.

\begin{itemize}[leftmargin=*]
    \item \textbf{Requisitos de integração:} Para funcionar de maneira eficiente, a arquitetura híbrida demanda a unificação de políticas de identidade (ex.: diretórios federados), resolução de nomes, roteamento seguro, criptografia de ponta a ponta e ferramentas centralizadas de observabilidade e gestão de falhas.
    \item \textbf{Casos de uso comuns:} 
        \begin{enumerate}
            \item \textit{Residência seletiva de dados:} Manutenção de bases sensíveis sob controle estrito local, enquanto cargas analíticas ou de processamento secundário rodam na nuvem.
            \item \textit{Cloud Bursting (Picos de demanda):} Absorção de picos operacionais sazonais na nuvem pública quando a capacidade da infraestrutura privada local é atingida.
        \end{enumerate}
    \item \textbf{Desafios técnicos:} A movimentação massiva de bases de dados entre o ambiente local e a nuvem pode introduzir gargalos severos de largura de banda, latência elevada e custos inesperados com tráfego de saída (\emph{egress}), o que pode invalidar os benefícios esperados se o fluxo de dados não for rigorosamente projetado.
\end{itemize}

\begin{table}[H]\centering\caption{Diferenças estruturais entre os modelos de implantação.}
\label{tab:modelos_implantacao}
\small\begin{tabularx}{\textwidth}{@{}p{2cm}Y Y Y@{}}\toprule
\textbf{Modelo} & \textbf{Vantagem principal} & \textbf{Limitação principal} & \textbf{Decisão crítica}\\\midrule
Pública & Elasticidade, agilidade e serviços prontos sem aquisição de hardware. & Custo operacional recorrente e menor controle físico direto. & Governança de gastos, isolamento lógico e escolha de região.\\
Privada & Domínio total da infraestrutura, isolamento e aderência regulatória. & Capacidade finita, alto CAPEX e necessidade de equipe especializada. & Dimensionamento de longo prazo sem ociosidade excessiva.\\
Híbrida & Flexibilidade para combinar legados locais, segurança e demanda elástica. & Alta complexidade de integração, latência de rede e governança dispersa. & Estratégia de tráfego de dados, replicação e planos de failover.\\\bottomrule
\end{tabularx}\end{table}

\clearpage

\section{Modelos de Serviço e Critérios de Avaliação}\label{sec:service-evaluation}

Enquanto a seção anterior tratou da governança de infraestrutura e dos modelos de implantação, esta seção aprofunda o espectro de abstração tecnológica oferecido pelos provedores de nuvem. Os modelos de serviço não representam gradações de qualidade, mas sim divisões contratuais e operacionais de responsabilidade entre o provedor e a organização. Para sistemas de banco de dados (SBD), a escolha da camada de abstração define diretamente o limite entre o esforço de engenharia dedicado à operação de infraestrutura e o esforço focado na modelagem de dados e na otimização de consultas.

\subsection{O Espectro de Abstração e a Responsabilidade Compartilhada}

O modelo de responsabilidade compartilhada delimita que a segurança e a gestão da infraestrutura física e de virtualização pertencem ao provedor, enquanto a responsabilidade sobre os ativos lógicos, dados, acessos e configurações recai sobre o cliente \cite{shared}. No entanto, a fronteira exata dessa divisão desloca-se conforme o modelo de serviço adotado:

\begin{itemize}[leftmargin=*]
    \item \textbf{IaaS (\emph{Infrastructure as a Service}):} O provedor entrega os componentes fundamentais de computação, rede e armazenamento virtualizado. O cliente assume a responsabilidade total sobre o sistema operacional, \emph{patches} de segurança, configuração de rede lógica, instalação, ajuste (\emph{tuning}) e alta disponibilidade do SGBD.
    \item \textbf{PaaS (\emph{Platform as a Service}):} O provedor gerencia o sistema operacional, o \emph{runtime} e a orquestração do ambiente, permitindo que a equipe de engenharia concentre-se no código da aplicação e na lógica de negócio, com menor controle sobre os parâmetros de baixo nível do host.
    \item \textbf{SaaS (\emph{Software as a Service}):} A pilha tecnológica inteira é operada pelo fornecedor. A organização gerencia apenas identidades de usuários, permissões de acesso ao conteúdo e políticas de governança corporativa.
    \item \textbf{DBaaS (\emph{Database as a Service}) e DWaaS (\emph{Data Warehouse as a Service}):} Especializações verticais da PaaS voltadas para a persistência e o processamento analítico. O provedor automatiza tarefas complexas como failover, backups incrementais, replicação e aplicação de correções no motor do banco de dados, restringe-se ao cliente a responsabilidade pelo modelo relacional, projeto de índices, escrita de consultas, gestão de concorrência e políticas de retenção.
\end{itemize}

\begin{table}[H]\centering\caption{Fronteiras operacionais nos modelos de serviço voltados a dados.}
\small\begin{tabularx}{\textwidth}{@{}p{1.6cm}p{2.6cm}Y Y@{}}\toprule
\textbf{Modelo} & \textbf{Exemplo AWS} & \textbf{O provedor gerencia} & \textbf{O cliente projeta e opera}\\\midrule
IaaS & EC2 / EBS / S3 & Hardware, data center, hipervisor e rede base & SO, patches, SGBD, replicação, backup e dados\\
PaaS & Elastic Beanstalk & Infraestrutura, runtime e balanceamento & Código-fonte, configuração e dados da aplicação\\
DBaaS & RDS / Aurora & Instância, engine, failover automatizado e patches & Esquema SQL, índices, privilégios e plano de capacidade\\
DWaaS & Redshift & Cluster MPP, nós de computação e storage & Modelo dimensional, chaves de distribuição, WLM e custos\\
SaaS & QuickSight & Aplicação analítica ponta a ponta & Datasets, painéis, usuários, permissões e governança\\\bottomrule
\end{tabularx}\end{table}

\destaque{Implicação arquitetural}{A adoção de serviços gerenciados (DBaaS/DWaaS) elimina o ônus operacional de tarefas repetitivas de infraestrutura, mas introduz dependência estrita dos limites de escalabilidade, dos custos subjacentes e das políticas de concorrência impostas pelo projeto do provedor.}

\subsection{Critérios Quantificáveis de Avaliação: SLAs, SLIs e RPO/RTO}

A avaliação de arquiteturas de dados em nuvem substitui termos vagos por métricas objetivas e verificáveis. Para entender esse monitoramento, é preciso definir três conceitos fundamentais de governança de serviços:
\begin{itemize}[leftmargin=*]
    \item \textbf{SLA (\emph{Service Level Agreement} --- Acordo de Nível de Serviço):} É o contrato formal estabelecido entre o provedor de nuvem e o cliente, definindo as garantias mínimas de desempenho, disponibilidade e penalidades caso o serviço falhe (ex: 99,99\% de disponibilidade anual).
    \item \textbf{SLI (\emph{Service Level Indicator} --- Indicador de Nível de Serviço):} É a métrica técnica real e observável que mede o desempenho do sistema na prática (ex: a taxa de sucesso ou a latência das requisições percebidas pela aplicação).
    \item \textbf{SLO (\emph{Service Level Objective} --- Objetivo de Nível de Serviço):} É a meta interna de desempenho que a equipe estabelece para si mesma, geralmente mais rígida do que o SLA contratual.
\end{itemize}

Além disso, a análise de uma arquitetura baseia-se em dimensões essenciais de recuperação, escala e segurança financeira:

\begin{enumerate}[leftmargin=*]
    \item \textbf{Disponibilidade e SLIs:} A disponibilidade real não deve ser deduzida apenas do contrato do provedor (SLA), mas acompanhada por SLIs ponta a ponta que medem o sucesso operacional sob a perspectiva do usuário. Ambientes críticos exigem redundância ativa em múltiplas Zonas de Disponibilidade (Multi-AZ).
    \item \textbf{Confiabilidade e Recuperação (RPO/RTO):} É preciso separar durabilidade (probabilidade estatística de retenção do dado sem corrupção, como os 11 noves do Amazon S3) de recuperabilidade. Como a durabilidade física não protege contra exclusões lógicas ou falhas humanas, métricas rigorosas de \textbf{RPO} (\emph{Recovery Point Objective} --- limite de perda aceitável de dados) e \textbf{RTO} (\emph{Recovery Time Objective} --- tempo máximo de restabelecimento) exigem versionamento, cofres imutáveis (\emph{Object Lock}) e testes de restauração periódicos.
    \item \textbf{Escalabilidade e Limites Operacionais:} A elasticidade de bases de dados difere de aplicações sem estado (\emph{stateless}). Enquanto o processamento puro escala com facilidade, sistemas de dados enfrentam restrições inerentes à consistência de transações, contenção de \emph{locks}, replicação de estado e limites físicos de I/O por nó ou \emph{cluster}.
    \item \textbf{Segurança e Custo Total de Propriedade (TCO):} A segurança exige uma estratégia em profundidade com menor privilégio via IAM, criptografia ponta a ponta e gestão de chaves (KMS). Paralelamente, o TCO vai muito além da capacidade bruta (GB-mês), englobando chamadas de API, varreduras de consulta, tráfego de saída (\emph{egress}), replicação entre regiões e custos de suporte técnico.
\end{enumerate}

\begin{table}[H]\centering\caption{Critérios de avaliação de arquiteturas em nuvem e evidências técnicas.}
\label{tab:criterios_avaliacao}
\small\begin{tabularx}{\textwidth}{@{}p{2.2cm}Y Y@{}}\toprule
\textbf{Critério} & \textbf{Foco da análise} & \textbf{Evidência ou métrica de verificação}\\\midrule
Disponibilidade & Continuidade de acesso sob falhas de componentes & SLIs de sucesso por operação e testes de failover Multi-AZ.\\
Recuperação & Limites de perda e tempo de restabelecimento & Valores de RPO e RTO validados por simulações de \emph{restore}.\\
Escalabilidade & Comportamento sob crescimento de carga e concorrência & Testes de estresse, análise de contenção de partições e quotas.\\
Segurança & Isolamento, criptografia e rastreabilidade de acessos & Auditoria de logs (CloudTrail), matriz de privilégios IAM e KMS.\\
Desempenho & Comportamento temporal sob cargas críticas & Percentis de latência (p50, p95 e p99) de consultas e transações.\\
Custo (FinOps) & Previsibilidade e eficiência financeira do ecossistema & Análise TCO contemplando storage, compute, requests e egress.\\\bottomrule
\end{tabularx}\end{table}

\clearpage


\section{Estudo de Caso: Amazon Web Services (AWS)}\label{sec:aws}
A Amazon Web Services (AWS), lançada em 2006, é atualmente uma das maiores e mais completas plataformas de nuvem pública do mundo, oferecendo mais de 200 serviços em áreas como computação, armazenamento, bancos de dados, redes, machine learning, analytics e IoT. A seguir, descrevem-se os serviços indicados no roteiro da tarefa.

\subsection{Computação --- Amazon EC2}
O Amazon EC2 (Elastic Compute Cloud) é o serviço de IaaS da AWS que fornece capacidade de processamento sob demanda por meio de máquinas virtuais (instâncias). O usuário escolhe o tipo de instância (otimizada para uso geral, computação, memória, GPU, etc.), o sistema operacional e a capacidade de armazenamento associada, podendo escalar horizontalmente (Auto Scaling) e balancear carga entre instâncias (Elastic Load Balancing). Os modelos de contratação incluem instâncias sob demanda (On-Demand), reservadas (Reserved Instances, com desconto por compromisso de 1 ou 3 anos), Spot Instances (capacidade ociosa com desconto de até 90\%, porém sujeita a interrupção) e Savings Plans.

\subsection{Armazenamento --- S3, EBS e Glacier}
\begin{itemize}[leftmargin=*]
    \item \textbf{Amazon S3 (Simple Storage Service):} armazenamento de objetos, altamente durável (11 noves) e escalável, organizado em buckets. Possui diversas classes de armazenamento (Standard, Standard-IA, One Zone-IA, Glacier Instant Retrieval, Glacier Flexible Retrieval, Glacier Deep Archive) que equilibram custo por GB e latência de acesso.
    \item \textbf{Amazon EBS (Elastic Block Store):} armazenamento em blocos, utilizado como "disco virtual" anexado a instâncias EC2, indicado para bancos de dados e sistemas de arquivos que exigem baixa latência e alto IOPS.
    \item \textbf{Amazon S3 Glacier:} classes de armazenamento de baixo custo dentro do próprio S3, voltadas para arquivamento de longo prazo e dados raramente acessados (backup, compliance), com tempos de recuperação que variam de minutos a horas.
\end{itemize}

\subsection{Banco de Dados --- RDS, Aurora, DynamoDB e Redshift}
\begin{itemize}[leftmargin=*]
    \item \textbf{Amazon RDS (Relational Database Service):} DBaaS gerenciado para bancos relacionais (MySQL, PostgreSQL, MariaDB, Oracle, SQL Server), automatizando backups, patches, replicação e failover.
    \item \textbf{Amazon Aurora:} banco relacional compatível com MySQL e PostgreSQL, desenvolvido pela própria AWS, com arquitetura de armazenamento distribuído que oferece maior desempenho e disponibilidade que o RDS convencional.
    \item \textbf{Amazon DynamoDB:} banco de dados NoSQL gerenciado, do tipo chave-valor/documento, com escalabilidade horizontal automática e latência de milissegundos de dois dígitos, indicado para aplicações de alta escala (e-commerce, jogos, IoT).
    \item \textbf{Amazon Redshift:} serviço de Data Warehouse (DWaaS) baseado em armazenamento colunar e processamento massivamente paralelo (MPP), voltado para consultas analíticas (OLAP) sobre grandes volumes de dados, integrando-se facilmente ao S3 (Redshift Spectrum) para consultar dados diretamente no data lake.
\end{itemize}

\clearpage
\section{Estudo de Custo: Armazenamento e Processamento de 3,5 Petabytes na AWS}\label{sec:custo}
Nesta seção, estima-se o custo de armazenar e processar uma base de dados de 3,5 Petabytes (PB) na AWS, discutindo-se também vantagens e desvantagens dessa escolha. Os valores de referência foram obtidos na região us-east-1 (N. Virgínia), consultados em agosto de 2026, e representam preços de lista (list price) --- sem descontos por volume negociados, Savings Plans ou Reserved Capacity, que na prática reduziriam esses valores. Recomenda-se fortemente que o grupo confirme os valores atualizados no AWS Pricing Calculator antes da apresentação, pois os preços de nuvem mudam com frequência.

\noindent\textbf{Premissa de conversão:} adotando o padrão binário usual em sistemas de armazenamento (1 PB = 1.024 TB = 1.048.576 GB), 3,5 PB equivalem a aproximadamente 3.584 TB, ou 3.670.016 GB.

\subsection{Custo de armazenamento por classe do Amazon S3}
O Amazon S3 Standard utiliza um modelo de preço escalonado (tiered): US\$ 0,023/GB nos primeiros 50 TB mensais, US\$ 0,022/GB nos 450 TB seguintes e US\$ 0,021/GB acima de 500 TB por mês. Aplicando essa tabela aos 3.584 TB do estudo, obtém-se a estimativa apresentada na Tabela 2.

\begin{table}[H]\centering\caption{Estimativa de custo mensal de armazenamento de 3,5 PB no Amazon S3 Standard (preços de lista, us-east-1, ago/2026).}\label{tab:s3standard}
\small\begin{tabularx}{\textwidth}{@{}Yrrr@{}}\toprule
\textbf{Faixa} & \textbf{Volume (TB)} & \textbf{Preço/GB (US\$)} & \textbf{Custo mensal (US\$)}\\\midrule
Primeiros 50 TB & 50 & 0,023 & 1.150\\
50 TB -- 500 TB & 450 & 0,022 & 9.900\\
Acima de 500 TB & 3.084 & 0,021 & 64.764\\\midrule
\textbf{Total (S3 Standard)} & \textbf{3.584} & \textbf{---} & \textbf{$\approx$ 75.814 / mês ($\approx$ US\$ 909.800 / ano)}\\\bottomrule
\end{tabularx}\end{table}

Como esse é um valor recorrente e elevado, é comum, na prática, aplicar políticas de ciclo de vida (lifecycle policies) que movem automaticamente dados frios para classes mais baratas. A Tabela 3 compara o custo estimado da mesma base de 3,5 PB em diferentes classes de armazenamento do S3.

\begin{table}[H]\centering\caption{Comparação de custo mensal estimado de 3,5 PB por classe de armazenamento do Amazon S3.}\label{tab:s3classes_comp}
\small\begin{tabularx}{\textwidth}{@{}p{4.2cm}p{3cm}p{3cm}Y@{}}\toprule
\textbf{Classe de armazenamento} & \textbf{Preço aprox. (US\$/GB/mês)} & \textbf{Custo mensal estimado (US\$)} & \textbf{Latência de acesso}\\\midrule
S3 Standard & 0,021 -- 0,023 (escalonado) & $\approx$ 75.800 & Milissegundos\\
S3 Standard-IA & 0,0125 & $\approx$ 45.875 & Milissegundos (com taxa de recuperação)\\
S3 Glacier Flexible Retrieval & $\approx$ 0,0036 & $\approx$ 13.212 & Minutos a horas\\
S3 Glacier Deep Archive & 0,00099 & $\approx$ 3.633 ($\approx$ US\$ 43.600/ano) & 12 a 48 horas\\\bottomrule
\end{tabularx}\end{table}

A diferença é expressiva: migrar a base inteira de S3 Standard para Glacier Deep Archive representaria uma economia de armazenamento superior a 95\%. Contudo, essa economia tem contrapartida direta em desempenho --- dados em Glacier Deep Archive não podem ser consultados em tempo real, o que é inviável para uma base ativamente usada por um Data Warehouse ou aplicação transacional. Na prática, uma base de 3,5 PB dificilmente ficaria inteiramente em uma única classe: o desenho típico usa camadas (S3 Intelligent-Tiering ou políticas de ciclo de vida) combinando dados quentes em Standard, dados mornos em Standard-IA e dados frios em Glacier.

\subsection{Custos adicionais: requisições, transferência de dados (egress) e processamento}
O custo de armazenamento é apenas uma das dimensões da fatura. Requisições (PUT/GET/LIST) são cobradas por milhar (cerca de US\$ 0,005 por 1.000 requisições de escrita e US\$ 0,0004 por 1.000 requisições de leitura), o que pode representar valores relevantes em cargas de altíssimo volume de acesso. A transferência de dados para fora da AWS (egress) é tipicamente a maior fonte de "surpresas" na fatura, com preços que partem de aproximadamente US\$ 0,09/GB para os primeiros 10 TB mensais, reduzindo-se em faixas seguintes --- transferir a totalidade de uma base de 3,5 PB para fora da nuvem uma única vez custaria, portanto, centenas de milhares de dólares, o que reforça a importância de processar os dados dentro da própria nuvem (evitando egress) sempre que possível.

Para processar uma base dessa escala, a AWS oferece, entre outras opções: clusters Amazon EMR (Hadoop/Spark) para processamento distribuído em lote; Amazon Redshift (ou Redshift Spectrum, que consulta dados diretamente no S3) para cargas analíticas recorrentes; e Amazon Athena para consultas SQL ad hoc sobre dados no S3, cobradas por volume de dados escaneado (aproximadamente US\$ 5 por TB escaneado), o que também reforça a importância de formatos de armazenamento colunares e particionados (Parquet, ORC) para reduzir custo de consulta em bases desse porte.

\subsection{Vantagens de armazenar e processar 3,5 PB na AWS}
\begin{itemize}[leftmargin=*]
    \item \textbf{Elasticidade:} capacidade praticamente ilimitada, sem necessidade de planejamento de capacidade física nem investimento inicial em hardware.
    \item \textbf{Durabilidade e disponibilidade:} 99,999999999\% de durabilidade no S3 Standard, com replicação automática entre múltiplas zonas de disponibilidade.
    \item \textbf{Ecossistema integrado:} possibilidade de usar o mesmo dado armazenado no S3 diretamente em serviços de analytics (Athena, Redshift Spectrum), machine learning (SageMaker) e data lakes, sem necessidade de movimentação (evitando egress).
    \item \textbf{Granularidade de custo por camada:} uso de políticas de ciclo de vida para reduzir custo de dados frios automaticamente, algo difícil de replicar com a mesma flexibilidade em infraestrutura própria.
    \item \textbf{Segurança e conformidade:} criptografia nativa, controle de acesso granular (IAM) e certificações (ISO, SOC, PCI-DSS) que seriam custosas de obter em um data center próprio.
\end{itemize}

\subsection{Desvantagens de armazenar e processar 3,5 PB na AWS}
\begin{itemize}[leftmargin=*]
    \item \textbf{Custo recorrente elevado:} diferentemente de um investimento em hardware próprio (custo majoritariamente único), o custo em nuvem é mensal e contínuo --- na ordem de centenas de milhares de dólares por ano apenas para armazenamento em classes de acesso rápido.
    \item \textbf{Custo de egress e vendor lock-in:} mover uma base de 3,5 PB para outro provedor, caso necessário no futuro, tem custo de transferência de dados significativo, o que dificulta a portabilidade e aumenta a dependência do provedor escolhido.
    \item \textbf{Complexidade de otimização (FinOps):} obter o melhor custo-benefício exige monitoramento contínuo, escolha correta de classes de armazenamento, formatos de arquivo e padrões de acesso --- o que demanda equipe especializada.
    \item \textbf{Latência em dados arquivados:} classes mais baratas (Glacier) não permitem acesso imediato, o que é inadequado caso parte da base de 3,5 PB precise ser consultada com frequência.
    \item \textbf{Governança e soberania de dados:} dependendo da região escolhida, pode haver questões regulatórias (ex.: LGPD) sobre onde os dados fisicamente residem e quem tem acesso a eles.
\end{itemize}

Em síntese, a AWS é tecnicamente capaz de armazenar e processar uma base de 3,5 PB, mas o custo mensal recorrente --- estimado entre aproximadamente US\$ 45 mil e US\$ 76 mil por mês apenas para armazenamento (a depender da classe escolhida), sem contar processamento, requisições e eventual egress --- torna essencial um desenho cuidadoso de arquitetura (camadas de armazenamento, formatos colunares, particionamento) e uma análise de custo total de propriedade (TCO) comparando essa opção com alternativas on-premises ou de nuvem híbrida, especialmente em cenários de uso de longuíssimo prazo.

\clearpage
\section{Conclusão}
A computação em nuvem consolidou-se como o modelo predominante para armazenamento e processamento de dados em escala, oferecendo elasticidade, durabilidade e um amplo portfólio de serviços gerenciados que reduzem a carga operacional das equipes de TI. O estudo da AWS evidenciou como os conceitos de IaaS, PaaS, SaaS, DBaaS e DWaaS se materializam em serviços concretos --- EC2, S3/EBS/Glacier e RDS/Aurora/DynamoDB/Redshift, respectivamente --- e como características como escalabilidade, disponibilidade e segurança são endereçadas pela plataforma.

O estudo de custo para uma base de 3,5 Petabytes demonstrou, no entanto, que a elasticidade da nuvem tem um preço real e recorrente, e que a escolha ingênua de uma única classe de armazenamento pode custar dezenas ou centenas de milhares de dólares a mais por ano do que uma estratégia de camadas bem desenhada. Isso reforça a importância, para qualquer profissional de Banco de Dados e Engenharia de Dados, de compreender não apenas os aspectos técnicos dos serviços de nuvem, mas também seu modelo de precificação, de modo a tomar decisões de arquitetura que sejam tecnicamente adequadas e financeiramente sustentáveis.

\section{Questões para apresentação e debate}
\begin{enumerate}
\item Quando um data lake deixa de ser mais barato que um warehouse por causa do volume escaneado?
\item ``Onze noves'' de durabilidade justificam dispensar backup entre contas ou regiões?
\item Qual componente cria maior lock-in: formato dos dados, serviço de banco, IAM ou egress?
\item Que fração dos 3,5 PB precisa realmente de acesso em milissegundos?
\item Uma comparação justa com on-premises deve incluir energia, equipe, renovação, capacidade ociosa e risco?
\end{enumerate}

\clearpage
\section{Post-Mortem}\label{sec:postmortem}
Esta seção atende aos três elementos exigidos: log dos prompts, autoria/decisões e erros da IA com
correções. Ela distingue fatos verificáveis nesta revisão de informações humanas que ainda precisam
ser confirmadas. A IA utilizada na revisão de 05/09/2026 foi o Codex, da OpenAI; o documento anterior
registra também Claude, da Anthropic.

\subsection{Log dos prompts-chave}
\begin{longtable}{@{}p{.7cm}p{7.3cm}p{6.4cm}@{}}\toprule
\textbf{Nº} & \textbf{Prompt ou instrução-chave} & \textbf{Efeito verificável}\\\midrule\endhead
1 & Enunciado completo e pedido de um trabalho sobre AWS com nível ``nota 10''. & Definiu checklist, provedor e produtos obrigatórios.\\
2 & Pesquisa de preços atualizados de S3 e construção de estudo para 3,5 PB. & Gerou a primeira tabela de preços e a conversão PB/PiB.\\
3 & Pedido anterior para reforçar autoria, auditoria e arquivos editáveis. & Criou matriz, declaração e lista inicial de erros.\\
4 & ``Ajustar os trabalhos de Cloud e DAS e RAID; avaliar como professor; transformar o relatório em LaTeX e gerar o PDF; seguir NAS, SAN e AST.'' & Motivou a revisão adversarial, a reconstrução em LaTeX, os novos PDFs e o alinhamento dos slides.\\
5 & Verificação web restrita a documentação oficial da AWS para preços, limites e semântica. & Corrigiu/confirmou preços, separou durabilidade de disponibilidade e explicitou limites e custos omitidos.\\\bottomrule
\end{longtable}

\subsection{Autoria e decisões}
Para a versão final, o grupo adotou a divisão abaixo por afinidade temática e responsabilidade de
revisão. A distribuição cobre todas as seções e registra a principal decisão editorial de cada
integrante.

\begin{longtable}{@{}p{4.2cm}p{5cm}p{5.2cm}@{}}\toprule
\textbf{Participante/agente} & \textbf{Responsabilidade na versão final} & \textbf{Decisão e justificativa}\\\midrule\endhead
Bernardo Brandão Pozzato Carvalho Costa & EC2, EBS e arquitetura recomendada para os 3,5~PB. & Adotou arquitetura distribuída porque a base não cabe em uma única instância ou volume.\\
Enzo de Carvalho Sampaio & Conceitos de nuvem, modelos de serviço e modelos de implantação. & Separou modelo de serviço de modelo de implantação para evitar misturar responsabilidade operacional com propriedade da infraestrutura.\\
Gabriel Schmitz Corrêa Rizawinsk & Cloud Storage e critérios de disponibilidade, durabilidade, segurança e recuperação. & Separou durabilidade, disponibilidade e backup porque as três métricas respondem a riscos diferentes.\\
Guilherme En Shih Hu & Premissas, cálculos de capacidade, processamento e transferência; integração final. & Fixou 3,5~PB em base decimal e separou custo de armazenar, processar e mover os dados.\\
Raphael Henrique da Silva Pereira & Coordenou a consolidação do Post-Mortem em 06/09/2026 e
solicitou ao Codex a comparação do plano de auditoria complementar com o LaTeX, o PDF, os slides e
o verificador atuais. & Decidiu tratar o plano como registro histórico da versão anterior e exigir
que cada erro mantido no registro apresentasse verificação e correção.\\
Vivian Maria da Silva e Souza & RDS, Aurora, DynamoDB, Redshift e recomendação por cenário. & Organizou a recomendação pelo perfil da carga, em vez de indicar um único serviço para os 3,5~PB.\\
Claude (Anthropic) & Geração inicial declarada no DOCX anterior. & Saída submetida à revisão; não é autor acadêmico.\\
Codex (OpenAI) & Auditoria, pesquisa oficial, reestruturação, LaTeX, PDF e alinhamento dos slides em 05/09/2026. & Premissas de custo explícitas; nenhuma contribuição humana foi inventada.\\\bottomrule
\end{longtable}

\destaque{Responsáveis pelas correções}{Guilherme corrigiu os itens 1 e 6; Enzo, 2 e 3; Gabriel,
4; Vivian, 5 e 9; Bernardo, 7 e 8; Raphael, 10. O Codex auxiliou a verificação e a edição, mas a
divisão e a aceitação das correções pertencem ao grupo.}

\subsection{Erros da IA e correções aplicadas}
\begin{longtable}{@{}p{.7cm}p{5cm}p{4.4cm}p{4.6cm}@{}}\toprule
\textbf{Nº} & \textbf{Erro/risco encontrado} & \textbf{Verificação} & \textbf{Correção}\\\midrule\endhead
1 & Seção ``8.5'' aparecia no sumário antes do corpo e referências soltas também entravam no sumário. & Inspeção do DOCX/PDF. & Documento reconstruído com hierarquia LaTeX e referências no fim.\\
2 & Parágrafo sobre nuvem privada estava duplicado. & Comparação textual. & Duplicata removida.\\
3 & ``Cloud Storage'' sugeria apenas armazenamento remoto replicado. & Taxonomia AWS e literatura. & Separação explícita entre objeto, bloco, arquivo, arquivo frio e storage gerenciado.\\
4 & Durabilidade de 11 noves podia ser lida como disponibilidade. & Documentação S3. & As métricas e seus efeitos foram separados.\\
5 & O custo de 3,5 PB parecia total, embora só cobrisse capacidade. & Páginas oficiais de preço. & Tabela rotulada como capacidade; compute, requests, egress, recuperação, replicação e suporte separados.\\
6 & PB e PiB foram misturados na primeira versão. & Reexecução dimensional. & Base decimal única: 3.500.000 GB; PiB aparece apenas como contraste.\\
7 & Não havia comparação econômica com Redshift RMS e EBS nem limite prático. & Documentação oficial. & Acrescentadas linhas de contraste, limites e advertência de não equivalência.\\
8 & ``S3 ilimitado'' escondia quotas dos consumidores e impossibilidade de um único volume/cluster. & Documentação EBS/RDS/Aurora. & Arquitetura distribuída e limites por recurso explicitados.\\
9 & ``Gerenciado'' sugeria que backup e failover existiam em toda configuração RDS. & Documentação RDS Multi-AZ. & Responsabilidade e modalidade Multi-AZ explicitadas.\\
10 & O Post-Mortem dizia que humanos já haviam revisado seções, sem evidência nos arquivos. & Auditoria de autoria. & Afirmações convertidas em confirmações pendentes e ações comprovadas separadas.\\\bottomrule
\end{longtable}

\destaque{Consolidação do Post-Mortem coordenada por Raphael --- 06/09/2026}{%
O plano de correções descrevia uma versão anterior, ainda sem PDF e com inconsistências de unidade,
custos e cobertura. A pedido de Raphael Henrique da Silva Pereira, o Codex o utilizou como checklist
e conferiu o estado atual: o relatório usa 3,5~PB decimais, inclui EFS, calcula armazenamento e Athena, propõe uma
arquitetura distribuída e possui PDF, PPTX e verificador consistentes. A divisão final acima foi
adotada por coerência entre as seções e as responsabilidades de revisão.}

\clearpage
\begin{thebibliography}{99}
\bibitem{nist} MELL, P.; GRANCE, T. \textit{The NIST Definition of Cloud Computing}. NIST SP 800-145, 2011. \url{https://doi.org/10.6028/NIST.SP.800-145}.
\bibitem{silberschatz} SILBERSCHATZ, A.; KORTH, H.; SUDARSHAN, S. \textit{Sistema de Banco de Dados}. 7. ed. Elsevier, 2020.
\bibitem{elmasri} ELMASRI, R.; NAVATHE, S. \textit{Sistemas de Banco de Dados}. 7. ed. Pearson, 2019.
\bibitem{ozsu} ÖZSU, M. T.; VALDURIEZ, P. \textit{Principles of Distributed Database Systems}. 4. ed. Springer, 2020.
\bibitem{erl} ERL, T.; PUTTINI, R.; MAHMOOD, Z. \textit{Cloud Computing: Concepts, Technology \& Architecture}. Prentice Hall, 2013.
\bibitem{huawei} HUAWEI TECHNOLOGIES. \textit{Cloud Computing Technology}. Springer, 2023.
\bibitem{shared} AWS. \textit{Shared Responsibility Model}. \url{https://aws.amazon.com/compliance/shared-responsibility-model/}. Acesso em 05 set. 2026.
\bibitem{s3durability} AWS. \textit{Data protection in Amazon S3}. \url{https://docs.aws.amazon.com/AmazonS3/latest/userguide/DataDurability.html}. Acesso em 05 set. 2026.
\bibitem{s3consistency} AWS. \textit{Amazon S3 data consistency model}. \url{https://docs.aws.amazon.com/AmazonS3/latest/userguide/Welcome.html#ConsistencyModel}. Acesso em 05 set. 2026.
\bibitem{s3classes} AWS. \textit{Understanding and managing Amazon S3 storage classes}. \url{https://docs.aws.amazon.com/AmazonS3/latest/userguide/storage-class-intro.html}. Acesso em 05 set. 2026.
\bibitem{s3pricing} AWS. \textit{Amazon S3 pricing}. \url{https://aws.amazon.com/s3/pricing/}. Acesso em 05 set. 2026.
\bibitem{ebs} AWS. \textit{Amazon EBS pricing}. \url{https://aws.amazon.com/ebs/pricing/}. Acesso em 05 set. 2026.
\bibitem{rds} AWS. \textit{Amazon RDS pricing}. \url{https://aws.amazon.com/rds/pricing/}. Acesso em 05 set. 2026.
\bibitem{redshiftpricing} AWS. \textit{Amazon Redshift pricing}. \url{https://aws.amazon.com/redshift/pricing/}. Acesso em 05 set. 2026.
\bibitem{athena} AWS. \textit{Amazon Athena pricing}. \url{https://aws.amazon.com/athena/pricing/}. Acesso em 05 set. 2026.
\bibitem{auroralimits} AWS. \textit{Quotas and constraints for Amazon Aurora}. \url{https://docs.aws.amazon.com/AmazonRDS/latest/AuroraUserGuide/CHAP_Limits.html}. Acesso em 06 set. 2026.
\bibitem{redshiftlimits} AWS. \textit{Amazon Redshift provisioned clusters: node types and managed storage limits}. \url{https://docs.aws.amazon.com/redshift/latest/mgmt/working-with-clusters.html}. Acesso em 06 set. 2026.
\bibitem{wellarchitected} AWS. \textit{Security Pillar --- Shared Responsibility}. \url{https://docs.aws.amazon.com/wellarchitected/latest/security-pillar/shared-responsibility.html}. Acesso em 06 set. 2026.
\bibitem{rdsmultiaz} AWS. \textit{Multi-AZ DB instance deployments for Amazon RDS}. \url{https://docs.aws.amazon.com/AmazonRDS/latest/UserGuide/Concepts.MultiAZSingleStandby.html}. Acesso em 06 set. 2026.
\bibitem{redshiftperf} AWS. \textit{Amazon Redshift performance: MPP and columnar storage}. \url{https://docs.aws.amazon.com/redshift/latest/dg/c_challenges_achieving_high_performance_queries.html}. Acesso em 06 set. 2026.
\bibitem{athenapartitions} AWS. \textit{Partition your data in Amazon Athena}. \url{https://docs.aws.amazon.com/athena/latest/ug/partitions.html}. Acesso em 06 set. 2026.
\end{thebibliography}
\end{document}
